\section{Introduction}\label{sec:intro}

A bank's allowance for credit losses is one number on its balance sheet, and
under the current expected credit loss standard (CECL) it is a forecast: the
losses the bank expects over the whole remaining life of its loans. The
standard lets the bank choose how to estimate it \citep{asc326}, the
supervisors' policy statement calls the estimate ``inherently imprecise''
\citep{ips2023acl}, and the SEC expects registrants to document ``the
identified range'' of expected losses and why they picked the estimate they
did within it \citep{sab119}. Nobody says how wide that range is. This paper
measures one part of it that is easy to overlook: the part that comes from
the credit model alone.

The problem is that many credit models fit the same history about equally
well. A lender that trains gradient-boosted trees, constrained trees and a
spline logistic regression on the same loans, and keeps every model whose
held-out accuracy is within a stated tolerance of the best, has several
defensible default models. Each implies an allowance. If they agree, the
choice does not matter. If they do not, the reserve depends on a decision
that accuracy alone does not settle, and the documented range is wider than
any single model shows.

I measure this on Freddie Mac's public loan-level data: \num{data.orig_rows}
fixed-rate mortgages and \num{data.perf_rows} billion loan-months of
performance. At each of \num{const.n_dates} year-ends from 2005 to 2025 I
train \num{const.n_models} default models on what a lender would have known
then, keep those within one bootstrap standard error of the best held-out log
loss, and compute the lifetime allowance of the same 200{,}000 portfolio
loans under every one, holding fixed everything else a CECL estimate needs: the
prepayment model, the economic forecast, the reversion to long-run averages
after the two-year forecast period, and the loss-severity table. The design
and six predictions were committed before any loan-level file was opened.

Four findings stand out.

\begin{enumerate}[leftmargin=*, itemsep=2pt]
\item \emph{Equally accurate models disagree substantially.} The range of the
allowance across the set has a median of \num{res.rel.median}\% of the most
accurate model's allowance, and up to \num{res.usd.max} billion dollars on
Freddie Mac's non-defaulted active book (December \num{res.usd.max_year}).
Stricter sets shrink it but do not remove it: among the models within two
standard errors of their paired difference from the best (a median of
\num{ph.paired2.size}) the range is a median \num{ph.paired2.rel}\%, and within
one paired standard error (a median of \num{res.paired1.size} models, a single
model at \num{p2.paired1.single} dates) it is \num{res.paired1.rel}\%. The
recorded range exceeds the range between one- and three-year reversion to
long-run averages at every date, by a median factor of
\num{res.ratio_policy.median}, although that is only part of the discretion
the standard grants over the forecast (\cref{sec:limits}).
\item \emph{Retraining the same model moves it at least as much.} The recorded
test, against bootstrap refits, holds at \num{H1.n} of \num{const.n_dates}
dates, but those refits are less accurate than the model they refit. Refits of
the most accurate configuration on independent random samples of the same
size, which are as accurate as the original fit to within a fraction of a
standard error, are the fairer benchmark (added after the results). At the
\num{ph.rs.disj.n} dates where those samples are disjoint, twenty such refits
span a median \num{p2.disj.rsm}\% of the allowance against \num{p2.disj.lfl}\%
for twenty randomly chosen equally accurate models (\cref{sec:refits}).
\item \emph{No range covers what happened.} Realized losses fall inside
the set's range at one of \num{H5.n_elig} dates with eight or more years of
outcomes, and inside the range of the refits at \num{p2.in_refit} (bootstrap)
and \num{p2.in_rsm} (independent samples). They were \num{res.miss2006} times the top of the range for the
December 2006 portfolio and below it at every date from 2010 to 2017. Model uncertainty and
forecast uncertainty are different things, and a range built from models says
nothing about the second.
\item \emph{IFRS 9 staging is fragile.} Across the set a median
\num{res.flip.median}\% of loans change stage under the twice-the-origination-PD
test. The recorded prediction that the set flips more loans than the refits
holds at \num{H6.n} dates, but the medians are almost equal, and in comparisons
added after the results twenty independent-sample refits flip more loans than
twenty set models at all but \num{ph.rs.flipmore} dates (\num{ph.rs.flip}\%
against \num{ph.H6.flipset}\%): the prediction holds, but not for the reason it
was made.
\end{enumerate}

\noindent Two predictions failed or held only narrowly. The relative range was
not unusually large on the eve of the crisis (H3 fails), and monotone
constraints narrowed it at only \num{H4.n} of \num{const.n_dates} dates.

The paper makes three contributions. It gives, to my knowledge, the first measurement on public
data of how far a CECL allowance moves across a formally defined set of
equally accurate loan-level models, with the other inputs held fixed. It
separates that model-choice range from training-sample noise, from the discretion over
reversion and from realized losses, over two decades that include a crisis
and a pandemic. And it supplies a reproducible procedure, released with the
code, that a lender, validator or auditor can run to produce the identified
range SAB 119 asks for.

\Cref{sec:background} sets out the accounting and supervisory rules and the
related work, \cref{sec:data,sec:method} the data and the design,
\cref{sec:results} the results and \cref{sec:discussion} what they mean for
preparers, auditors and model validators.
